% Candidate introduction for research-question option 1; not integrated into the report.
\section{Introduction}
Automatic chord recognition produces a timed harmonic summary of music audio. In a major/minor vocabulary, each chord label combines a root with a quality. Recovering these labels from a recording is complicated by overlapping harmonics, unequal note strengths and changes within an analysis window. Chroma provides a compact representation by folding spectral evidence into twelve pitch classes, but it also removes register information. Chord models then interpret this evidence, while temporal decoding favours consistency across frames\ \cite{mcvicar2014,pauwels2019,mueller2021fmp}.

These stages interact. Smoothing can remove brief prediction errors but also suppress short, genuine chords. Cho et al.\ \cite{cho2010} found that the effective transition penalty depended on the acoustic model and its score separation. A penalty shared across representations therefore need not produce comparable smoothing. Vocabulary choices further affect how predictions match the reference labels\ \cite{humphrey2015}. These considerations motivate examining changes in error type alongside overall accuracy.

In our preliminary guitar experiments, direct note projection improved root accuracy over long-window chroma, while same-root major/minor errors increased after shared decoding. The corresponding raw error totals were similar, but this does not establish that their locations or score margins were alike. Our research question is: \textbf{Does projection's increase in same-root quality errors persist after decoder calibration?} We will compare raw predictions, a shared switching penalty and penalties selected separately on development data. Paired error analysis and controlled transitions will measure the remaining quality-error gap and the detection of genuine changes.
